\documentclass[10pt,conference]{IEEEtran}
\usepackage[T1]{fontenc}
\usepackage{url}
\usepackage{booktabs}
\usepackage{graphicx}

\title{An Inspectable Rule-First Transpiler for Legacy-Language Subsets: A Tool Demonstration}

% ICSE 2027 Tool Demonstration review is single-anonymous: replace these fields.
\author{\IEEEauthorblockN{[Author Name]}\\
\IEEEauthorblockA{[Department and Institution]\\
[City, Country]\\
[Contact Email]}}

\begin{document}
\maketitle

\begin{abstract}
Legacy-code translation tools can produce plausible target programs while obscuring
which source constructs were handled and what was actually checked. We demonstrate a
compiler-design prototype that translates declared subsets of COBOL, FORTRAN, Pascal,
and line-numbered BASIC into JavaScript or Python, and a project-specific Python-like
subset into JavaScript. Its web interface exposes tokenization, normalized intermediate
data, deterministic output, unsupported constructs, optional complete-program model
proposals, and finite validation results. The model output remains an unaccepted
candidate. On a 20-case task-grouped pilot, both targets passed 17 supported cases and
matched all 17 separately authored expected outputs; three COBOL loop cases were
correctly reported unsupported. These expected outputs have not been checked against
native source compilers. A separate refresh at pushed commit ade68dc confirmed prompt
v3 and the configured Gemini model but returned no candidate after three provider
attempts per target. The prototype therefore demonstrates inspectable boundaries and
failure reporting, not an LLM effectiveness gain or general semantic equivalence.
Demonstration video: \url{[PUBLIC_YOUTUBE_VIDEO_URL]}
\end{abstract}

\section{Introduction}
Modernizing legacy code requires more than producing syntactically plausible output:
users need to know which constructs were translated, where behavior is uncertain, and
what evidence supports a candidate. Prior translation systems combine code generation
with execution, testing, or explicit feature analysis in their evaluated settings
\cite{yang2024,zhang2025}. Our contribution is narrower: a runnable teaching/research
prototype whose interface makes its declared subset and validation boundary visible.

The prototype does not establish a new translation algorithm, full language coverage,
or an effectiveness advantage. It provides (1) custom subset front ends and a shared
normalized representation, (2) deterministic rule-first translation, (3) optional
complete-program LLM proposals for unsupported input, and (4) finite target execution
and separately authored output checks. It is intended for researchers and students
inspecting translation decisions, not production migration teams.

\section{Tool and Demonstration}
The front ends recognize documented subsets of COBOL, FORTRAN, Pascal, BASIC, and a
project-specific Python-like language. The first four languages can target JavaScript or
Python; the Python-like demonstration path targets JavaScript. This is not full-dialect
support. The normalized representation and deterministic code generator expose the
rule-based path. Parser-reported unsupported input can trigger an optional request for
a complete target program. The application retains deterministic partial output,
displays the model candidate separately, reports candidate syntax/runtime status, and
never marks the candidate semantically accepted.

\begin{figure}[t]
\centering
\fbox{\begin{minipage}{0.96\columnwidth}
\centering
Source subset $\rightarrow$ parser and normalized representation $\rightarrow$
deterministic target code $\rightarrow$ target checks\\[2pt]
Unsupported construct $\rightarrow$ optional full-program candidate
$\rightarrow$ separate review status
\end{minipage}}
\caption{Implemented path. Candidate generation is optional; it is not an automatic
acceptance or repair loop.}
\label{fig:pipeline}
\end{figure}

The demonstration walks through a supported loop in a legacy source language, shows
its intermediate representation and target output, then enters a COBOL \texttt{PERFORM}
loop outside the declared subset. The UI shows the unsupported statements, retained
deterministic output, fallback status, and any candidate checks. The benchmark panel
shows per-case results and distinguishes expected-unsupported recognition from
translation success. The fallback prompt treats the complete source as untrusted text.

JavaScript candidates are syntax-checked and executed using Node's permission mode,
with a timeout and a reduced child environment. Python candidates are syntax-checked
only and are not executed. Neither behavior is a complete security sandbox or semantic
acceptance gate. The current [candidate acceptance policy](CANDIDATE_ACCEPTANCE.md)
requires multi-case evidence and human review before any future acceptance workflow.

\section{Evaluation and Limitations}
We ran the deterministic benchmark against JavaScript and Python. The test split has
20 records in two whole-task groups: maximum-of-two-numbers and for-loop tasks. It
contains 17 expected-translatable records and three expected-unsupported COBOL
\texttt{PERFORM} records. The 17 supported target executions matched their separately
authored expected outputs in both targets (17/17); all three unsupported cases were
recognized (3/3). End-to-end status was 17/20 for each target because expected-unsupported
cases are reported as \texttt{not\_run}, not translation passes (Table~\ref{tab:results}).

\begin{table}[t]
\centering
\caption{Deterministic held-out pilot.}
\label{tab:results}
\begin{tabular}{lrr}
\toprule
Measure & JavaScript & Python \\
\midrule
Test cases & 20 & 20 \\
Supported outputs matching expected output & 17/17 & 17/17 \\
Expected-unsupported cases recognized & 3/3 & 3/3 \\
End-to-end translation passes & 17/20 & 17/20 \\
\bottomrule
\end{tabular}
\end{table}

The expected outputs are project-authored from fixed constants and loop ranges and
match a separate reference evaluator (20/20). This checker does not import the
transpiler or its source interpreter; it is not a native source compiler or a second
human review. The corpus has only two test
task groups and is not a representative estimate of translation accuracy. The original
24-case paired refresh at pushed commit \texttt{ade68dc} used
\texttt{gemini-3.1-flash-lite} and \texttt{legacy-modernizer-llm-v3}. The provider
returned no candidate after three request attempts for the unsupported case in each
target, and the hybrid success rate did not exceed the rule-only rate. This run supports
no LLM effectiveness claim.

Contracts, risk scores, failure fingerprints, and learned-rule examples are prototype
metadata. No controlled ablation has shown that they improve correctness or review
time. We make no scalability or multilingual-bias claim. Python candidate execution,
independent source-language oracles, a larger task-grouped corpus, and repeated
controlled model trials remain future work. The complete run provenance and case JSON
are linked in the repository's \texttt{data/EVIDENCE\_V3\_REFRESH.md}.

\section{Availability}
The source and usage instructions are publicly available at
\url{https://github.com/yoshi-hue/transpiler}. The tool runs locally using Python,
FastAPI, and Uvicorn; the README documents installation, configuration, and execution.
The optional model provider must be configured by the user. No API credential is
included in the artifact.

\begin{thebibliography}{2}
\bibitem{yang2024}
Z. Yang et al., ``Exploring and Unleashing the Power of Large Language Models in
Automated Code Translation,'' in \emph{Proc. ACM Joint European Software Engineering
Conference and Symposium on the Foundations of Software Engineering (ESEC/FSE)}, 2024,
doi: 10.1145/3660778.

\bibitem{zhang2025}
H. Zhang, C. David, M. Wang, B. Paulsen, and D. Kroening, ``Scalable, Validated Code
Translation of Entire Projects using Large Language Models,'' \emph{Proceedings of the
ACM on Programming Languages}, vol. 9, PLDI, pp. 1616--1641, 2025,
doi: 10.1145/3729315.
\end{thebibliography}
\end{document}
